\section{MANEJO Y USO DE ALARMAS SONORAS Y GEIGER MÜLLER}

\textbf{OBJETIVO:} El alumno debe darse cuenta de que estos equipos son la mejor herramienta para desempeñar su trabajo dentro de ciertos parámetros de seguridad y que, dadas las características de las radiaciones, es la única manera de saber cuando estamos dentro de un campo de altos niveles de rapidez de exposición.

\textbf{EQUIPO:}
\begin{itemize}
    \item Equipo portátil Detector de radiaciones con Alarma Sonora.
    \item Equipo portátil Detector de Radiaciones tipo Geiger-Müller.
    \item Contenedor de Gamagrafía.
    \begin{itemize}
        \item Conteniendo una fuente radiactiva con la suficiente actividad como para poder tener en sus costados niveles de rapidez de exposición de alrededor de 50 a 100 mR/h.
    \end{itemize}
\end{itemize}

\textbf{DESARROLLO:}
\begin{enumerate}
    \item \textbf{PARA LA ALARMA SONORA}
        \begin{itemize}
            \item Encender su indicador de Alarma Sonora, colocando el "switch" en posición de encendido y en la escala que se quiera, aproximarlo a casi contacto del contenedor; se escuchará la indicación audible (sonido intermitente), de forma clara y fuerte. En caso de no escuchar dicha señal, proceder a reemplazar la batería verificando que su colocación sea correcta. Encender nuevamente y aproximarla al contenedor. Si no funciona, notificarlo de inmediato a su instructor.
        \end{itemize}
    \item \textbf{PARA EL EQUIPO DETECTOR DE RADIACIONES TIPO GEIGER-MÜLLER}
        \begin{itemize}
            \item Encender el equipo y poner la perilla selectora de escala en posición ``batería'', observar que la aguja se desplace a la posición que se indica en la carátula del equipo. Si la aguja no llega a la posición adecuada, cambiar las baterías por nuevas vigilando que la instalación sea correcta. Colocar el detector a 50 cm.
        \end{itemize}
\end{enumerate}